\documentclass[11pt]{article}
\usepackage{amsmath,amssymb,amsthm}
\newtheorem{theorem}{Theorem}\newtheorem{proposition}{Proposition}
\title{Paid delay: guarded source recursion without silent initialization}
\date{16 November 2026}
\begin{document}\maketitle
\paragraph{Audit.} Independently reviewed addendum to the birth-stable initial-row and random-initializer feedback constructions. This changes the timing interface by moving an initial draw into one forced paid event. It is a functor on native hom categories with a natural independent-evaluation comparison, not a strong, lax or oplax preservation claim for horizontal units/composition. The initial witness and every actual coefficient remain present. The earlier well-founded silent theorem is not weakened to accept cyclic initialization.

\section{Raw initial histories are not decoded states}
At a designated input $a$ of an initial-row arena, the raw horizon-zero history set is $Z_a$, carrying its probabilities and actual coefficients. Its decoded state is $\iota_a(z)$, but a control table may inspect the full $z$. A later history retains this seed followed by all selected scopes and their raw expanded outcomes, including zero and off-policy outcomes. If tests beginning at arbitrary retained configurations are also part of a controller interface, those tests remain separately tagged; no equivalence of that larger interface is inferred just from designated initial executions.

For example, two distinct fair initial witnesses may select the same configuration with two independent fair-bit primitive rows. An admissible paid-scope controller can choose the diagonal coupling at seed zero and the antidiagonal coupling at seed one. Collapsing the initial history to its decoded state destroys this distinction although all primitive marginals agree. Thus preserving the seed in the following timing translation is material, not optional bookkeeping.

\section{The actual paid-delay functor}
For a birth-stable initial-row arena $P:A\rightsquigarrow B$ with initial row
$(Z_a,\nu_a,L_{a,z})$ and initial-state map $\iota_a$, define $D(P)$ as follows.
Add a separate boot configuration $\bot_a$ over every input $a$; keep ALL
original configurations, including unreachable ones. A boot configuration
has no emitted port and has exactly one enabled paid event, with a fresh
boot tag distinct from original event identifiers. Its raw row is exactly
$(Z_a,\nu_a,L_{a,z})$, and outcome $z$ selects the original state $\iota_a(z)$.
The new initial row is a singleton of weight one and coefficient $I$, selecting
$\bot_a$. Original configurations, paid rows, transitions, port data, full
grade $C$ and terminal bodies are unchanged. A boot state is not terminal.
The boot event has no coenabled distinct event, so birth stability there is
vacuous; original birth-stability equations are unchanged.

For a native cell $H:P\to Q$, send $\bot_a$ to $\bot_a$, use its original
configuration/local maps elsewhere, and use its actual initial-row arrow as
the boot-event row arrow. The new singleton initializer map is forced. The
old initial square is exactly the boot transition square. All other native
squares and persistence laws are the original ones; the full grade map is
unchanged.
\begin{theorem}
These definitions give a functor $D$ on each native initial-row arena hom
category, and a natural actual native isomorphism
\[
 \operatorname{Eval}_{\rm ind}(D(P))\cong
 \operatorname{Eval}_{\rm ind}(P).
\]
The comparison is identity on the full content grade and preserves the raw
initial witness and all actual context coefficients.
\end{theorem}
\begin{proof}
The cell formulas compose exactly: the boot row map is the composite of the
original initial-row maps, including their augmented coefficient arrows.
Identity is equally immediate on all retained configurations and full grades.
A complete designated execution of $D(P)$ has its forced singleton initial
witness, one boot outcome $z$, and a complete original paid trace. Delete only
the forced singleton marker and reinterpret $z$ as the original initial
witness. This is a raw bijection, including zero outcomes, and preserves the
probability $\nu_a(z)$ times the original paid law. Its actual coefficient
map is the unitor
$I\otimes L_{a,z}\otimes K_{\rm trace}\cong L_{a,z}\otimes K_{\rm trace}$.
Nothing discards $L_{a,z}$. Terminal full bodies agree. The same initial and
paid native arrows act on both descriptions, proving naturality.
\end{proof}

Designated control histories after the boot retain its outcome $z$ in the
raw past. Therefore a supplied original policy on designated seed histories
can be implemented after the boot without forgetting seed-dependent choices;
the boot scope itself is forced and its single-row law fixed. This does not
assert an equivalence of all possible policy structures on every unreachable
configuration or every newly introduced boot-state test.

\section{What delay does not preserve}
Let $1$ be the structural zero-event unit. Its initial state is terminal,
whereas $D(1)$ initially has one enabled event. No native cell exists in either
direction, since enabled-event bijections would compare zero with one.
Likewise $D(1;1)$ has one paid boot, whereas $D(1);D(1)$ has two sequential
paid boots. Their initial singleton raw rows force any initial comparison;
after the first event one side is terminal and the other is not. Hence no
native comparison exists in either direction. The natural evaluated
isomorphism is not a native operational contraction and cannot make $D$
strong, lax or oplax on horizontal units/composition in this cell class.

Nor is fullness asserted. For a concrete failure, let $Q$ have one terminal
state $t$ per its sole input, one unary output, singleton initial row selecting
$t$, and unit coefficients. Let $P$ contain that state and initializer plus
an unreachable state $x$ with no ports and one deterministic event leading
to $t$. Give both the same full unary unit grade. There is no native $P\to Q$
map because $x$ has one enabled event and every $Q$ state has none. But
$D(P)\to D(Q)$ can send $x$ to the new boot state of $D(Q)$, its event to the
boot event, $t$ to $t$, and the designated boot to the designated boot.
All row, transition and terminal-body squares hold. Thus additional native
maps arise from the extra boot states; they are not images of old maps.

\paragraph{Coefficient-doctrine pin.}
Throughout this addendum $\operatorname{Rec}_U(F)$ means the terminating
invocation-tree recurrence with the TENSOR of the original context
coefficients at its invocation nodes. The ``old atomic realization'' is the
one-paid-event tensor-coefficient realization, not the earlier
product-coefficient all-coupling family. For the corresponding same-law
product-coefficient process there is only the directed comparison
$\lambda:\bigotimes_v E_v\to\prod_v^{\mathcal E}E_v$, obtained at component
$v$ by augmenting every other factor; raw witnesses, weights and the full
content grade are unchanged. This is the comparison proved in
\emph{Tensor Causal Execution and Native Comparison}
(\texttt{\detokenize{Tensor_Causal_Execution_and_Native_Comparison.tex}}).
It need not be invertible: for two augmented copies of $k^2$ with first-coordinate
augmentation, it maps a four-dimensional tensor to a three-dimensional
augmentation pullback and kills the tensor of the two augmentation-kernel
vectors. No tensor/product identification or recovery of every correlated
member of the larger all-coupling family is claimed.

\section{Guarded feedback and exact original-source recovery}
Every initializer of $D(P)$ selects a port-free boot state. Its all-raw
initial-loop relation is empty. Consequently the random well-founded
feedback compiler applies to $D(P):A+U\rightsquigarrow B+U$ for ANY original
birth-stable $P$ of those endpoints, even when $P$'s original initial-loop
relation has cycles. The compiled arena still retains the boot outcome and
coefficient at each invocation. In particular a zero-probability seed is
not omitted to obtain this guard.

\begin{theorem}[Guarded recursive recovery]
Let $F$ be an original finite normalized weighted-polynomial process
$A+U\to B+U$, and $J(F)$ its strong zero-paid initial-row realization.
Then $D J(F)$ is natively isomorphic to the old one-paid-event atomic
realization of $F$, regarded in the birth-stable initial-row doctrine by
its singleton initializer. Feedback therefore preserves the old atomic
paid-control interface and its actual independent recursive semantics:
\[
 \operatorname{Eval}_{\rm ind}(\operatorname{Fb}_U(DJ(F)))
 \cong \operatorname{Rec}_U(F)
\]
on finite terminating raw invocation trees, with the full $R_U C$ and
actual context/content coefficients. This statement does not require the
undelayed initializer relation of $J(F)$ to be well founded.
\end{theorem}
\begin{proof}
The configurations of $D J(F)$ are a port-free boot over each input and one
terminal state per original raw context witness. Its boot row is exactly
the original normalized context row, with its actual augmented coefficient.
Its terminal shape and output labels are the original body. These are exactly
the old atomic realization's data, including every zero witness and full
unused grade. Local event/port maps identify the singleton boot tag and
ordered output positions, and all actual row/coefficient maps are identities.
Both singleton initial rows agree. This is a native isomorphism on ALL
configurations, not merely an evaluated equality.

Apply the guarded feedback functor and transport complete paid-scope policies
through the induced initial/event raw bijections. The original atomic source
has exactly these configurations and scopes, so this policy conservation is
specific and justified; it is not the unqualified all-policy equivalence for
arbitrary $P$ excluded above. Independent invocation normal form groups one
boot event at each completed invocation. The raw boot outcome is precisely
that invocation's original $F$ witness. Its probability and actual coefficient
are unchanged, and the full terminal body is the same formal $R_U C$ tree.
Thus finite terminating trees and their actual native data agree. Nontermination
is permitted in the paid dynamics and contributes no terminal raw witness;
no silent initializer cascade or infinite coefficient tensor has been used.
\end{proof}

For general $P$, the same independent normal-form theorem and the natural
comparison $\operatorname{Eval}_{\rm ind}D(P)\cong\operatorname{Eval}_{\rm ind}P$
identify the terminating feedback with invocation recurrence of the latter
process, on its stated countable/subprobability interface. At an invocation,
paid cost is now one plus its original paid-event cost. This change of clock
is explicit. For $J(F)$ every completed invocation costs one, so the event
clock counts invocation nodes, not macro generations. Under countable endpoint
hypotheses, the cited cost-clock matching trace recovers the same terminating
process, without claiming equality of differently indexed unclosed clocks.

A fair geometric source that either returns or loops is the discriminating
example: its undelayed zero-paid $J(F)$ is excluded from finite all-raw silent
initialization, whereas $D J(F)$ is guarded and supports arbitrarily long
paid recursive runs. This is a timing/interface change that restores the
old guarded source mechanism, not a proof that the original silent guard
was unnecessary or that arbitrary correlated boot laws have been added.
\end{document}
